\section{Results and Discussion}

\subsection{From tensor transport to spectroscopic use}

The main value of the present framework is not only conceptual. Once quartic
and sextic centrifugal distortion constants are lifted to their tensor
representatives, the same machinery may be used in practical spectroscopic
workflows to:
\begin{enumerate}
\item compare parameter sets obtained in different reductions or
representations;
\item transport fitted or computed constants from one representation to
another;
\item identify numerically ill-conditioned representation changes before a new
fit is attempted;
\item monitor the first breakdown of the standard linear sector through the
diagnostics $H_{22}$ and, on the sextic side, the term linear in
$\tau(H_{22})$.
\end{enumerate}

This section is therefore organized around real use cases rather than around
formal properties alone. Dimethylsulfoxide is used to illustrate
representation instability within the standard sector. Formaldehyde and
thioformaldehyde illustrate the transport of sextic standard constants through
their canonical tensor content. The HBN isomer family is then used as a
medium-sized quartic test bed, already large enough to make the practical
choice of computational hierarchy relevant. Water and formaldehyde are then
used to show how the quartic diagnostic $H_{22}$ behaves when the standard
VPT2 picture is respectively under stress or still reliable. Taken together,
these examples show that the present framework is not merely a reformulation
of known theory, but a practical tool for spectroscopic analysis. The analyses
discussed below are implemented in the app workflow: the standard quartic and
sextic transport steps start directly from centrifugal distortion constants,
whereas the post-standard diagnostics $H_{22}$ and $\tau(H_{22})$ use
optional harmonic input (`.fchk` or `xyz + Hessian`).

From the validation point of view, the role of these examples is also
well-defined. DMSO validates quartic transport in a case where representation
conditioning is known to be problematic. H$_2$CO and H$_2$CS validate the
standard sextic transport machinery against experimental and variational
parameter sets. H$_2$O versus H$_2$CO validates the first-breakdown diagnostic
$H_{22}$ by showing that it grows precisely where the standard VPT2 picture is
known to be under greater strain. The HBN family finally shows that the same
logic remains meaningful on systems of intermediate size, where computational
hierarchy and transport strategy become practical rather than purely formal
issues.

The quantum-chemical data entering these applications were generated within
the Pisa Composite Scheme (PCS) framework. In particular, two members of the
PCS family play complementary roles. The hybrid PCS2 scheme (HPCS2), based on
the B3LYP/6-31+G* level of theory, provides reliable anharmonic contributions
to rotational spectroscopic parameters. The double-hybrid PCS3 scheme
(DPCS3), based on the revDSD-PBEP86-D3BJ functional combined with the 3F12-
basis set, provides a more accurate harmonic baseline through improved
equilibrium structures and harmonic force fields.\cite{pcs1234_2024,DPCS3_2023,BDPCS3_2023}

This level of accuracy is already spectroscopically meaningful. In the
DPCS3//HPCS2 strategy, supplemented when needed by simple bond-length
corrections as in the BDPCS3 protocol, relative errors on rotational
constants may be pushed to the order of $10^{-3}$, while quartic and sextic
centrifugal-distortion constants can be described at the few-percent level.
\cite{acr_2026,Boussessi_quartic,Boussessi_sextic,hormones_2024,fast_dvib_2025}
Equally important, these calculations are no longer restricted to textbook
benchmarks: workflows of this type are now feasible for molecules with several
tens of atoms, up to the $\sim 50$-atom scale.\cite{hormones_2024,fast_dvib_2025}
This is precisely why the transformation problem addressed here has become
important in practice. Once quartic and sextic constants can be computed with
this level of accuracy on medium-sized systems, a representation analysis is
needed not only for conceptual clarity but also for the practical comparison
between fits, direct calculations, and transported parameter sets.

This distinction matters directly for the present paper. The standard
tensor-transport machinery and the quartic $H_{22}$ diagnostic are anchored in
the harmonic model, whereas the standard sextic anharmonic sector depends on
the cubic force field. For medium-sized molecules this computational hierarchy
is now realistic: modern workflows already provide accurate quadratic force
fields and semi-diagonal cubic force constants on systems far beyond the
traditional small-molecule benchmark set.\cite{cubic-vib_2024,cubic-internal_2024,fast_dvib_2025,hormones_2024,RotBio_2023,AnnRevPC_2023}
In that regime, a low-cost
diagnostic of the limits of standard perturbative transport is especially
useful, because the harmonic information is already affordable and accurate
even when a full post-standard treatment is not.

Operationally, this leads to a natural three-level hierarchy on the sextic
side:
\begin{enumerate}
\item harmonic model only: geometry sector and first H$_{22}$-induced sextic
diagnostic;
\item harmonic model plus semi-diagonal cubic sector: reduced sextic
prediction at nearly linear cost;
\item full cubic field: genuine three-index remainder and complete standard
sextic prediction.
\end{enumerate}
This hierarchy is important for the present paper because it connects the
formal tensor decomposition directly to the computational resources that are
realistically available for medium-sized molecules.

The corresponding practical workflow is therefore straightforward.
Starting from fitted Watson constants alone, one may perform standard
quartic and sextic transport between representations and evaluate the
conditioning diagnostics associated with the lift--transport--reproject
procedure. If harmonic input is also available in the form of geometry
plus Cartesian Hessian (or equivalently an input \texttt{.fchk} file),
one obtains in addition the standard quartic tensor, the quartic
$H_{22}$ diagnostic, the sextic geometry sector, and the first sextic
diagnostic induced by $\tau(H_{22})$. If semi-diagonal cubic data are
added, the reduced sextic prediction becomes available at nearly linear
cost in the number of modes, whereas adding the full cubic field yields
the genuine three-index remainder and thus the complete standard sextic
prediction. In this sense, the present framework does not require a
single all-or-nothing computational protocol, but naturally supports a
hierarchy of inputs, costs, and outputs.

\subsection{DMSO: representation transport and numerical instability}

The choice of Watson reduction and principal-axis representation is known to
affect the numerical stability of centrifugal-distortion fits. This is
particularly clear for dimethylsulfoxide (DMSO), which was analyzed in detail
by Margul\`es \emph{et al.}\cite{Margules2010} and already served as a key case
in the previous version of this work.

DMSO is especially useful here because the rotational constants $A$ and $B$
are close in magnitude. As a consequence, representation changes involving the
$I^r \leftrightarrow III^r$ conversion may become poorly conditioned, even
though the underlying tensor content remains the same. In this situation the
present tensorial strategy is particularly informative. One starts from a
given quartic Watson parameter set, reconstructs the reduced quartic tensor by
pseudoinverse lifting, transports that tensor under the required axis
permutation, and finally reprojects the result onto the target Watson
representation. In the present work, this entire workflow was executed through
the app, exactly in the same form in which it would be used in practical
spectroscopic analysis.

When this is done for DMSO, the individual Watson constants can vary strongly
between representations, especially in the anisotropic sector. This is not a
failure of the tensor algorithm. It is precisely the phenomenon that the
framework is meant to disentangle. The large variation is driven by the
conditioning of the chosen parametrization, not by a change in the underlying
quartic rovibrational content. This is seen most clearly by comparing the
reconstructed reduced quartic tensor $\tau'$ rather than the raw Watson
constants: its spectral invariants remain unchanged, up to the expected
permutation similarity, across representation transport.

The DMSO example therefore plays two roles. First, it validates the tensor
transport procedure on a realistic spectroscopic case for which representation
dependence is known to be severe. Second, it shows why a tensor-level analysis
is operationally preferable to a direct comparison of Watson constants.
Quantities such as $s_{111}$ and the condition number of the transport provide
practical indicators of whether a nominally equivalent representation is
actually a stable one for fitting purposes.

The numerical content of Table~\ref{tab:DMSO_constants} illustrates this point
directly. In the $A$ reduction, the change from $I^r$ to $III^r$ leaves some
constants only moderately affected, but others vary dramatically. The clearest
example is $\delta_K$, which changes from a value of order $1$ kHz in the
$I^r$ representation to a value of order $10$--$50$ kHz in the $III^r$
representation. The corresponding $s_{111}$ values show the same trend: while
they remain very small in the stable representations, they grow sharply in the
ill-conditioned one. In the $S$ reduction the same phenomenon is visible in a
less extreme but still clear form, with the anisotropic constants responding
much more strongly than the isotropic ones to the representation change.

These numerical variations are too large to be interpreted as simple changes
in underlying rovibrational physics. They are instead the signature of a
representation problem: the same physical quartic content is being described
through a parametrization whose conditioning deteriorates when the transformation
denominators become small. This is exactly the kind of situation in which the
tensor formulation is useful, because it separates the physical content from
the instability of the Watson coordinates.

\begin{table}[h]
\centering
\caption{Quartic centrifugal distortion constants and corresponding
$s_{111}$ parameter for dimethylsulfoxide in different Watson
representations (kHz). Rotational constants used in the analysis are
$A=7036.6$, $B=6910.8$, $C=4218.78$ MHz for the experimental fit and
$A=6749.88$, $B=6716.88$, $C=4052.83$ MHz for the HPCS2 calculation.}

\begin{tabular}{lccccc}
\hline
Constant & $I^r$ (QM) & $I^r$ (fit) & $III^r$ (QM) & $III^r$ (fit) & $III^r$ (fit, transf) \\
\hline

\multicolumn{6}{c}{A reduction} \\

$\Delta_J$    & 3.9912 & 4.0509 & 6.5456 & 6.6268 & 6.6327 \\
$\Delta_{JK}$ & -4.4009 & -4.4528 & -12.0641 & -12.1574 & -12.1981 \\
$\Delta_K$    & 6.3961 & 6.7065 & 6.3961 & 6.6726 & 6.7065 \\
$\delta_J$    & 1.5568 & 1.4549 & -0.2796 & -0.1648 & -0.1640 \\
$\delta_K$    & 1.1788 & 1.2856 & -29.5443 & -46.8755 & -29.7583 \\
$s_{111}$     & $5.6\times10^{-6}$ & $5.3\times10^{-7}$ & $-4.86\times10^{-4}$ & $-1.3\times10^{-6}$ & $-1.73\times10^{-3}$ \\

\hline

\multicolumn{6}{c}{S reduction} \\

$D_J$         & 4.5954 & 3.4631 & 6.2204 & 6.0892 & 6.6843 \\
$D_{JK}$      & -8.0260 & 0.9259 & -12.9010 & -8.9372 & -8.7376 \\
$D_K$         & 9.4171 & 3.7676 & 9.4171 & 3.9893 & 3.7676 \\
$d_1$         & 1.5568 & 1.4549 & -0.0507 & 0.1641 & 1.5850 \\
$d_2$         & 0.3021 & 0.2939 & -0.1832 & -0.2717 & -0.1126 \\
$s_{111}$     & $7.8\times10^{-8}$ & $7.5\times10^{-8}$ & $3.6\times10^{-8}$ & $3.5\times10^{-8}$ & $3.5\times10^{-8}$ \\

\hline
\end{tabular}

\label{tab:DMSO_constants}
\end{table}

The same conclusion is seen more clearly at the tensor level. While the Watson
constants vary strongly, the spectral content of the reconstructed quartic
tensor is stable under representation transport. This is illustrated in
Table~\ref{tab:DMSO_tauprime}, where the eigenvalues and polynomial invariants
of the reduced quartic tensor are reported for the fitted and transported DMSO
parameter sets. The transported $III^r$ row reproduces the $I^r$ tensor
invariants within numerical precision, as expected for a single underlying
quartic tensor seen in different representations.

\begin{table}[h]
\centering
\caption{Spectral invariants of the reduced quartic tensor $\tau'$ for
dimethylsulfoxide obtained from different Watson representations. The
$III^r$ (fit, transf) row corresponds to the tensor obtained by transporting
the $I^r$ fit through the tensor algorithm. Equality of the eigenvalues and
invariants demonstrates the representation independence of the tensor
description. Small deviations in the independently fitted $III^r$ row arise
from rounding of the published spectroscopic constants.}

\begin{tabular}{lcccccc}
\hline
Representation & $\lambda_1$ & $\lambda_2$ & $\lambda_3$ & $I_1^{(\tau')}$ & $I_2^{(\tau')}$ & $I_3^{(\tau')}$ \\
\hline
$I^r$ (fit)           & -2.4752 & -8.7629 & -46.3875 & -57.6256 & 542.9974 & -1006.1386 \\
$III^r$ (fit)         & -2.4849 & -8.6332 & -46.4642 & -57.5824 & 538.0495 & -996.7901 \\
$III^r$ (fit, transf) & -2.4752 & -8.7629 & -46.3875 & -57.6256 & 542.9974 & -1006.1386 \\
\hline
\end{tabular}

\label{tab:DMSO_tauprime}
\end{table}

The numerical pattern in Table~\ref{tab:DMSO_tauprime} is the key counterpoint
to Table~\ref{tab:DMSO_constants}. The transported $III^r$ row reproduces the
$I^r$ eigenvalues and invariants essentially exactly, confirming that the
tensor algorithm preserves the intrinsic quartic content under representation
change. By contrast, the independently fitted $III^r$ row shows only small
deviations, entirely consistent with the finite precision and rounding of the
published constants. The important point is therefore not that the Watson
constants look different, but that the reconstructed tensor remains the same
within the expected numerical accuracy. This is the operational meaning of
representation-equivalent transport in a real spectroscopic case.

\subsection{H\textsubscript{2}CO and H\textsubscript{2}CS: standard sextics as transported tensor content}

The sextic side of the theory is illustrated most cleanly by small asymmetric
tops for which complete sextic parameter sets are available from both
experiment and variational calculations. Formaldehyde and thioformaldehyde are
natural examples because Ref.~\citenum{Dinu2022} reports sextic constants in
Watson's $S$ reduction and $I^r$ representation for both molecules.

In the present framework these sextic constants are not treated as seven
independent spectroscopic coordinates. Instead, they are decomposed as
\begin{equation}
\mathbf H_S = \mathbf B_{I^r}\,\boldsymbol\sigma + \mathbf r,
\end{equation}
where $\boldsymbol\sigma$ contains the five physical sextic coordinates of the
canonical sector and $\mathbf r$ is the residual in the complementary
two-dimensional space. The vector $\boldsymbol\sigma$ is the transported
physical content, whereas $\mathbf r$ acts as a diagnostic of how closely the
input parameter set follows the modeled standard sextic subspace.

For H$_2$CO and H$_2$CS, the canonical coordinates obtained from experiment and
from the VCI/PFIT parameter sets are very similar. This is exactly the type of
behaviour one wants from a transported tensor representative: even when the
individual Watson constants differ, the physical content extracted from them
remains stable. The residual norms are not exactly zero, but this is expected
for finite-precision spectroscopic parameter sets and should not be read as a
failure of the method. Rather, the residual provides a compact measure of the
distance between a real fitted parameter set and the idealized standard sextic
subspace used in the tensor model.

This example is important for the main thesis of the paper. It shows in a real
spectroscopic setting that the standard sextic sector can be handled by the
same lift--transport--reprojection logic as the quartic one, even though its
effective-Hamiltonian parametrization is richer. In other words, the sextic
problem does not require a distinct transport philosophy. It still belongs to
the same standard linear representation-equivalent sector.

At the same time, the sextic side of the present work is not merely a
transport exercise. The equations are written explicitly in terms of the
geometric quantities $c_1$, $c_2$, and $\tau$, plus a separate cubic-force
contribution. This makes it possible, at least in principle, to decompose the
standard sextic sector into geometry, semi-diagonal cubic terms, and the
genuinely three-index cubic remainder. That separation is especially relevant
for medium-sized molecules, where semi-diagonal and full cubic data may be
available at different levels of theory.

The numbers support this conclusion directly. For H$_2$CO, the experimental
and VCI/PFIT values of the canonical coordinates differ only slightly across
all five components: for example, $\sigma_3$ changes from about $3.26$ to
$3.36$ kHz and $\sigma_5$ from about $2.49$ to $2.56$ kHz. For H$_2$CS the
same pattern is seen on a larger absolute scale, with $\sigma_3$ near
$4.81$--$4.86$ kHz and $\sigma_5$ near $3.67$--$3.71$ kHz. Thus the tensor
coordinates retain the chemically expected trend, namely a systematically
larger sextic distortion for the sulfur analogue, while remaining stable under
the experiment/VCI comparison.

Operationally, this sextic decomposition and transport were again obtained
with the app, showing that the same tool handles both quartic and sextic
standard sectors in a uniform way.

The residual norms carry complementary information. They are small enough to
confirm that the published sextic constants lie close to the modeled standard
subspace, but not so small as to suggest an exact algebraic identity at the
level of rounded spectroscopic parameters. This is precisely the regime in
which the residual becomes a useful spectroscopic diagnostic rather than a
formal remainder with no interpretative value.

\subsection{HBN isomers: a medium-sized quartic test bed}

The present framework is not intended only for traditional benchmark
molecules. A useful intermediate test bed is provided by the HBN isomer
family, where accurate quartic constants, representation transforms, and
high-level PCS calculations are all available while the size of the systems is
already representative of modern rotational spectroscopy on medium-sized
molecules.

For the four HBN isomers (\emph{o}-HBN, \emph{p}-HBN, \emph{cis}-m-HBN, and
\emph{trans}-m-HBN), Table~\ref{table:hbn_main} collects the ground-state
rotational constants and quartic centrifugal distortion constants in the
$I^r$ representation and $A$ reduction, together with the transformed
constants in $III^r$ ($A$ reduction) and $I^r$ ($S$ reduction). The same table
also reports the gauge diagnostic $\log_{10}|s_{111}|$. These are exactly the
quantities that a working spectroscopist would want to inspect before deciding
whether a representation change is harmless, numerically delicate, or better
avoided altogether.

The numerical pattern is instructive. In the native $I^r$ ($A$) description,
the quartic constants are small and chemically well behaved across the four
isomers, and the corresponding $\log_{10}|s_{111}|$ values remain in the
$-10$ to $-11$ range. This is the fingerprint of a standard sector that is
still well conditioned. After transport to $III^r$ ($A$), however, the
anisotropic constants become much larger and the magnitude of $s_{111}$ grows
by roughly one to two orders, especially for the more asymmetric members of
the family. This reproduces, on a chemically broader series, the same lesson
already seen for DMSO: what changes strongly under representation transport is
the Watson parametrization, not necessarily the underlying quartic tensor.

The HBN family is especially useful because it bridges the gap between the
classical small-molecule benchmarks and the medium-sized systems that motivate
the present paper. The PCS calculations used here are representative of the
quality of data that is becoming available today for molecules with several
tens of atoms. In this regime, the tensor transport itself is already directly
useful, while the same systems also provide a realistic platform for testing
how far the sextic hierarchy can be pushed with harmonic input, semi-diagonal
cubic data, and eventually the genuine three-index cubic remainder.

\begin{table*}[ht]
\centering
\scriptsize
\setlength{\tabcolsep}{4pt}

\caption{Ground-state rotational constants (MHz) and quartic centrifugal
distortion constants (kHz) for the four HBN isomers in the $I^r$
representation (A reduction). Ground-state constants were obtained by
subtracting the HPCS2 vibrational corrections from the equilibrium values
computed at different levels. The diagnostic parameter $s_{111}$ is reported
as $\log_{10}|s_{111}|$. The table further includes quartic constants
transported to $III^r$ (A reduction) and to $I^r$ (S reduction) through the
tensor-transform workflow implemented in the app.}

\begin{tabular}{lccc ccc ccc ccc}
\toprule
& \multicolumn{3}{c}{\textbf{cis-$o$-HBN}}
& \multicolumn{3}{c}{\textbf{$p$-HBN}}
& \multicolumn{3}{c}{\textbf{cis-$m$-HBN}}
& \multicolumn{3}{c}{\textbf{trans-$m$-HBN}} \\
\cmidrule(lr){2-4}
\cmidrule(lr){5-7}
\cmidrule(lr){8-10}
\cmidrule(lr){11-13}
Parameter
& Exp. & HPCS2 & DPCS3
& Exp. & HPCS2 & DPCS3
& Exp. & HPCS2 & DPCS3
& Exp. & HPCS2 & DPCS3 \\
\midrule
\multicolumn{13}{c}{\textbf{Rotational constants (MHz)}} \\
$B_a$ & 3053.7 & 3037.2 & 3062.1 & 5613.0 & 5596.5 & 5638.8 & 3408.9 & 3398.4 & 3418.6 & 3403.1 & 3390.4 & 3412.8 \\
$B_b$ & 1511.3 & 1506.4 & 1512.6 & 990.5 & 986.5 & 990.7 & 1205.8 & 1201.0 & 1206.6 & 1208.5 & 1204.1 & 1209.2 \\
$B_c$ & 1010.8 & 1006.8 & 1012.3 & 841.9 & 838.7 & 842.7 & 890.7 & 887.3 & 891.8 & 891.8 & 888.4 & 892.9 \\
\midrule
\multicolumn{13}{c}{\textbf{Quartic constants in $I^r$ (A reduction) (kHz)}} \\
$\Delta_J$ & 0.038 & 0.0368 & 0.0370 & 0.020 & 0.015 & 0.016 & 0.041 & 0.0387 & 0.0390 & 0.040 & 0.0380 & 0.0383 \\
$\Delta_{JK}$ & 0.624 & 0.619 & 0.6204 & 0.231 & 0.221 & 0.223 & 0.027 & 0.0246 & 0.0250 & 0.020 & 0.0352 & 0.0356 \\
$\Delta_K$ & -0.100 & 0.381 & 0.384 & 9 & 0.918 & 0.922 & 1.180 & 1.173 & 1.1778 & 1.130 & 1.158 & 1.161 \\
$\delta_J$ & 0.011 & 0.010 & 0.010 & 0.005 & 0.003 & 0.003 & 0.014 & 0.014 & 0.014 & 0.014 & 0.014 & 0.014 \\
$\delta_K$ & 0.280 & 0.332 & 0.335 & 0.400 & 0.182 & 0.184 & 0.160 & 0.182 & 0.185 & 0.220 & 0.182 & 0.185 \\
$\log_{10}|s_{111}|$ & -10.43 & -10.33 & -10.33 & -11.27 & -11.71 & -11.71 & -11.36 & -11.20 & -11.19 & -11.00 & -11.19 & -11.18 \\
\midrule
\multicolumn{13}{c}{\textbf{$III^r$ (A reduction)}} \\
$\Delta_J$ & 0.5310 & 0.6411 & 0.6445 & 1.9072 & 0.4028 & 0.4068 & 0.3718 & 0.3818 & 0.3851 & 0.3990 & 0.3839 & 0.3869 \\
$\Delta_{JK}$ & -0.9650 & -1.2412 & -1.2499 & -4.0638 & -0.8502 & -0.8572 & -0.8222 & -0.8700 & -0.8788 & -0.9420 & -0.8703 & -0.8785 \\
$\Delta_K$ & 0.4500 & 0.6168 & 0.6223 & 2.1667 & 0.4563 & 0.4603 & 0.4633 & 0.4988 & 0.5046 & 0.5550 & 0.4963 & 0.5018 \\
$\delta_J$ & 0.1255 & 0.2450 & 0.2461 & 2.3052 & 0.2832 & 0.2847 & 0.2947 & 0.2924 & 0.2937 & 0.2805 & 0.2913 & 0.2921 \\
$\delta_K$ & 0.1150 & 0.2613 & 0.2635 & 2.4500 & 0.3205 & 0.3225 & 0.3750 & 0.3843 & 0.3870 & 0.3925 & 0.3805 & 0.3828 \\
\midrule
\multicolumn{13}{c}{\textbf{$I^r$ (S reduction)}} \\
$D_J$ & 0.0380 & 0.0368 & 0.0370 & 0.0200 & 0.0150 & 0.0160 & 0.0410 & 0.0387 & 0.0390 & 0.0400 & 0.0380 & 0.0383 \\
$D_{JK}$ & 0.6240 & 0.6190 & 0.6204 & 0.2310 & 0.2210 & 0.2230 & 0.0270 & 0.0246 & 0.0250 & 0.0200 & 0.0352 & 0.0356 \\
$D_K$ & -0.1000 & 0.3810 & 0.3840 & 9.0000 & 0.9180 & 0.9220 & 1.1800 & 1.1730 & 1.1778 & 1.1300 & 1.1580 & 1.1610 \\
$d_1$ & 0.5820 & 0.6840 & 0.6900 & 0.8100 & 0.3700 & 0.3740 & 0.3480 & 0.3920 & 0.3980 & 0.4680 & 0.3920 & 0.3980 \\
$d_2$ & -0.5600 & -0.6640 & -0.6700 & -0.8000 & -0.3640 & -0.3680 & -0.3200 & -0.3640 & -0.3700 & -0.4400 & -0.3640 & -0.3700 \\
\bottomrule
\end{tabular}

\label{table:hbn_main}
\end{table*}

The tensor perspective becomes clearer when one moves from transported Watson
constants to the representation-invariant content of the reconstructed quartic
tensor. Table~\ref{table:hbn_tauprime} reports the eigenvalues and polynomial
invariants of $\tau'$ for the HPCS2 HBN series. The two meta isomers are
especially informative: although their transported Watson constants differ
noticeably representation by representation, the tensor invariants show that
their intrinsic quartic content is in fact very similar. This is precisely
the kind of chemically meaningful comparison that is obscured at the level of
the raw Watson coordinates and recovered by the tensor lift.

\begin{table}[h]
\centering
\caption{Spectral invariants of the reduced quartic tensor $\tau'$ for the
HBN isomers reconstructed from the HPCS2 quartic constants.}

\begin{tabular}{lcccccc}
\toprule
Isomer & $\lambda_1$ & $\lambda_2$ & $\lambda_3$ & $I_1^{(\tau')}$ & $I_2^{(\tau')}$ & $I_3^{(\tau')}$ \\
\midrule
cis-$o$-HBN   & 0.9585 & -0.2946 & -5.1055 & -4.4416 & -3.6721 & 1.4417 \\
$p$-HBN       & 0.0267 & -0.0328 & -4.7299 & -4.7360 & 0.0280 & 0.0042 \\
cis-$m$-HBN   & -0.0243 & -0.2578 & -4.9727 & -5.2548 & 1.4090 & -0.0311 \\
trans-$m$-HBN & -0.0214 & -0.2508 & -4.9566 & -5.2288 & 1.3546 & -0.0266 \\
\bottomrule
\end{tabular}

\label{table:hbn_tauprime}
\end{table}

For the present paper, the importance of the HBN example is twofold. First, it
shows that the tensor transport machinery remains operational on molecules
that are no longer merely textbook benchmarks. Second, it makes the practical
motivation of the sextic reformulation more concrete: these are exactly the
systems for which the separation between harmonic, semi-diagonal cubic, and
genuine three-index cubic contributions becomes both scientifically useful and
computationally meaningful.

\subsection{Uridine: a demanding dual-level spectroscopic case}

The same framework can be pushed well beyond the medium-sized benchmark level.
Gas-phase uridine provides a demanding example because it combines a
29-atom molecular structure, high-resolution experimental data, and a
computational setting in which a dual-level treatment is both realistic and
informative.\cite{uridine_2015} In the present case, equilibrium rotational
constants are available at the DPCS3 and BDPCS3 levels, while vibrational
corrections and centrifugal distortion constants were obtained at the HPCS2
level. This makes uridine a natural example of the kind of practical
equilibrium-plus-correction strategy that motivates the present tensor
formulation.

In current spectroscopic work on large molecules, HPCS2-quality predictions of
rotational and centrifugal-distortion constants are already a realistic
standard. The present application goes one step further by combining the
improved BDPCS3 equilibrium structure with HPCS2 vibrational corrections,
which we denote as `BDPCS3//HPCS2`. This level of theory is still feasible on
workstations of the kind now routinely used for systems in the
30--50 atom range, but it yields rotational constants with an accuracy that is
already comparable to what much more demanding methods provide for smaller
molecules. On this basis, a meaningful analysis of centrifugal distortion
constants also becomes possible.

Table~\ref{tab:uridine_rotational_duallevel} compares the experimental
ground-state rotational constants with the theoretical `HPCS2//HPCS2`,
`DPCS3//HPCS2`, and `BDPCS3//HPCS2` values. The HPCS2 vibrational corrections
are clearly not negligible, especially for the largest rotational constant,
and their inclusion is essential for any meaningful comparison with
experiment. The dual-level comparison is also selective: the
`BDPCS3//HPCS2` combination gives the best agreement with the observed
ground-state constants, reaching a mean absolute percentage error of about
$0.20\%$, compared with about $0.44\%$ for `HPCS2//HPCS2` and about
$0.53\%$ for `DPCS3//HPCS2`. This is already a quantitatively significant
result for a 29-atom system.

\begin{table}[h]
\centering
\scriptsize
\caption{Uridine as a demanding dual-level case. Experimental ground-state
rotational constants are from Ref.~\citenum{uridine_2015}. The notations
`HPCS2//HPCS2`, `DPCS3//HPCS2`, and `BDPCS3//HPCS2` indicate that the
equilibrium constants are taken at the first level of theory and corrected by
HPCS2 vibrational contributions. The last row reports the mean absolute
percentage error (MAE\%) with respect to experiment.}

\begin{tabular}{lcccc}
\hline
Constant & Exp. ($B_0$) & HPCS2//HPCS2 & DPCS3//HPCS2 & BDPCS3//HPCS2 \\
\hline
$A$ (MHz) & 885.98961(14) & 881.93503 & 880.81087 & 883.90684 \\
$B$ (MHz) & 335.59622(35) & 334.44150 & 334.13848 & 335.10385 \\
$C$ (MHz) & 270.11210(20) & 268.70421 & 268.76154 & 269.53137 \\
\hline
MAE (\%) & --- & 0.441 & 0.531 & 0.199 \\
\hline
\end{tabular}
\label{tab:uridine_rotational_duallevel}
\end{table}

The quartic constants are equally informative. The supplied geometry is
already in the $I^r$ representation, so the HPCS2 quartic constants can be
transported directly to alternative representations through the tensor
workflow. Table~\ref{tab:uridine_quartic_comparison} reports the native $I^r$
and transported $III^r$ values. Only one quartic constant, $D_J$, was
resolved experimentally, with a value of $0.0115(10)$ kHz. The native $I^r$
prediction, $D_J = 0.00873$ kHz, lies in the same range, whereas the
transported $III^r$ set shows a much stronger reshuffling of the anisotropic
quartic content.

This point is spectroscopically useful. For a molecule of this size a full
quartic fit is difficult, and the experiment could resolve only one quartic
parameter. The calculated constants, however, indicate that at least three
quartic terms are likely to be spectroscopically relevant: besides $D_J$, the
native $I^r$ values of $D_{JK}$ and $D_K$ are also appreciable on the scale
of the fit. In this sense the tensor-transported theoretical quartics can
guide the construction of an experimentally stable Hamiltonian by indicating
which additional distortion terms are worth introducing before the fit becomes
overparameterized. The same calculations also indicate that the sextic
constants are far too small to be expected in the present experiment, so the
quartic sector remains the relevant target for any improved fit.

\begin{table}[h]
\centering
\scriptsize
\caption{Uridine quartic centrifugal distortion constants from the HPCS2
vibrational treatment, reported in the native $I^r$ representation and after
transport to $III^r$. The supplied geometry is already in the $I^r$
representation. Only $D_J$ was resolved experimentally.}

\begin{tabular}{lccc}
\hline
Quantity & Experiment & HPCS2 ($I^r$) & HPCS2 ($III^r$) \\
\hline
$\Delta_J$ (kHz) & --- & 0.010458 & 0.034881 \\
$\Delta_{JK}$ (kHz) & --- & 0.147760 & 0.074491 \\
$\Delta_K$ (kHz) & --- & -0.097394 & -0.097394 \\
$\delta_J$ (kHz) & --- & 0.000760 & -0.012972 \\
$\delta_K$ (kHz) & --- & 0.061064 & -0.050023 \\
\hline
$D_J$ (kHz) & 0.0115(10) & 0.008729 & 0.066718 \\
$D_{JK}$ (kHz) & --- & 0.158135 & -0.116532 \\
$D_K$ (kHz) & --- & -0.106040 & 0.061792 \\
$d_1$ (kHz) & --- & -0.000760 & 0.012972 \\
$d_2$ (kHz) & --- & -0.000865 & 0.015919 \\
\hline
\end{tabular}
\label{tab:uridine_quartic_comparison}
\end{table}

The contrast between representations and reductions is already visible at the
level of Table~\ref{tab:uridine_quartic_comparison}. In the native $I^r$
representation the symmetrically reduced constants identify a relatively clear
hierarchy, with $D_J$, $D_{JK}$, and $D_K$ standing out as the only quartic
terms of potentially observable size. After transport to $III^r$, by
contrast, the anisotropic quartic content is redistributed much more strongly:
the $A$- and $S$-reduced constants remain mathematically equivalent, but they
no longer provide the same intuitive or numerically stable basis for fitting.
This is precisely the type of situation in which tensor transport becomes
useful. A fit based on a single quartic constant can reproduce the data, but
it need not return the parameter that would be obtained in a more complete
representation-aware treatment. The comparison with theory therefore suggests
that a constrained fit, with the additional quartic terms introduced together
with physically motivated uncertainty bounds, would be more robust than a fit
in which the whole quartic sector is effectively collapsed onto a single
observed parameter.

Uridine therefore strengthens one of the practical conclusions of the present
paper. For small systems one may aim at benchmark-level predictions of all
spectroscopic constants. For larger systems, the more realistic target is a
selective but physically informed prediction. The present tensor framework is
already useful in that demanding regime: it supports accurate dual-level
rotational constants, transport of quartic constants between representations,
and a theory-guided indication of which distortion terms are worth trying in
an experimental fit.

\subsection{H\textsubscript{2}O versus H\textsubscript{2}CO: the quartic diagnostic \texorpdfstring{$H_{22}$}{H22}}

The standard linear sector is not the whole story, and the key question is how
to detect its first failure in a form that remains simple enough to be
practically useful. This is the role assigned here to the quartic channel
$H_{22}$.

Water and formaldehyde provide a particularly instructive comparison. In
water, the differences between standard VPT-based quartic constants and
higher-level variational results are known to be substantial. In
formaldehyde, the same discrepancies are much smaller. The present diagnostic
follows that trend.

For H$_2$O, the projected quartic contribution generated by $H_{22}$ is
clearly non-negligible on the scale of the fitted quartic constants and of the
VPT/VCI differences, although it is not by itself sufficient to explain the
largest discrepancies. For H$_2$CO, the same diagnostic is dramatically
smaller, by roughly two orders of magnitude. This is exactly the behaviour
expected from a useful first-breakdown indicator: it grows where the standard
picture is under strain and remains small where the standard picture is
already adequate.

The point is not that $H_{22}$ provides the whole correction beyond the
standard quartic sector. The point is that it is the first genuinely new
tensorial contribution, it is simple enough to isolate explicitly, and it can
be computed from the same basic rovibrational information already needed for
the standard problem. In this sense $H_{22}$ is not only a theoretical
channel; it is an operational diagnostic of when one should stop trusting the
standard linear transport picture as quantitatively sufficient. The key issue
is not that $H_{22}$ must be the largest post-standard quartic contribution in
absolute value, but that it is the first term expected to break the linear
equivalence of Watson parameter sets across representations. In this sense two
fits in different representations remain physically equivalent within the
present framework only as long as the $H_{22}$ diagnostic is negligible.

This comparison is important also from the software point of view: $H_{22}$ is
not introduced here as an external post-processing quantity, but as a
diagnostic computed within the same app workflow used for the standard tensor
transport once harmonic input is supplied.

The numerical comparison is most informative in relative rather than absolute
terms. For H$_2$O, the projected $H_{22}$ contributions correspond to
fractions at the $10^{-3}$--$10^{-1}$ level of the fitted quartic constants,
depending on the Watson parameter considered, and they are of the same order
as the smaller VCI--VPT differences even though they do not saturate the
largest ones. For H$_2$CO, the same diagnostic becomes smaller by about two
orders of magnitude and correspondingly negligible as a fraction of the
quartic constants. This relative collapse is the important point: it shows
that the first breakdown of linear transport is much less relevant in the
molecule for which the standard VPT2 picture is already known to be more
reliable.

\subsection{First sextic analogue: linear response in \texorpdfstring{$\tau(H_{22})$}{tau(H22)}}

The same idea extends naturally to the sextic side. Once the standard sextic
geometry functional is written in terms of the standard quartic tensor, the
first post-standard sextic contribution is obtained by replacing the latter by
\begin{equation}
\tau = \tau_{\mathrm{std}} + \tau(H_{22})
\end{equation}
and extracting the term linear in $\tau(H_{22})$.

For H$_2$O this sextic diagnostic is nonzero and measurable, with its largest
component at the level of about $10^3$ Hz. This is much smaller than the
dominant standard sextic constants, so it is not a candidate for a full
post-standard sextic theory by itself. But that is not the point. The
quantity is useful precisely because it is simple, low-cost, and sensitive to
the first way in which the new quartic tensorial level propagates into the
sestic problem.

The natural expectation is that this quantity should remain smaller in systems
for which standard VPT-based and higher-level results are already close, just
as the quartic $H_{22}$ diagnostic does for H$_2$CO. The full H$_2$CO sextic
numerical test requires the corresponding cubic data and may therefore be
added later without changing the logic of the present paper. At the current
stage, the role of the sextic analogue is already clear: it is the first
diagnostic of how the quartic breakdown feeds into the sextic geometry sector.

As for the quartic case, this quantity is available directly in the app
workflow once harmonic input is supplied, so that standard transports and
first-breakdown diagnostics can be inspected within a single operational
environment.

Here again the numerical scale matters. For water the largest component of the
linear sextic response is about $1.3\times 10^3$ Hz, with the dominant effect
appearing in the $\Phi_{aaa}$ channel. Relative to the dominant standard
sextic constants, this corresponds to a correction at about the $10^{-3}$
level: still perturbatively small, but clearly above numerical noise and
therefore defining a genuine post-standard scale. The quantity is therefore
well placed to act as a first warning sign of how the quartic breakdown feeds
into the sextic geometry sector.

\subsection{Scalability and modern spectroscopic scope}

One further point deserves to be stressed. The present framework is not only
suited to small benchmark molecules. Modern rotational spectroscopy and modern
quantum-chemical workflows now routinely reach systems with several tens of
atoms, including flexible organic molecules and biologically relevant species.
The tensor-lifting and transport logic developed here is entirely compatible
with that scale.

This matters for the positioning of the paper. If the framework were useful
only for small benchmark systems, it would remain largely methodological. In
fact, its intended use is broader: the same machinery can organize fitted and
computed centrifugal-distortion constants for medium-sized and large
molecules, help choose numerically stable reductions and representations
before refitting, and signal when the standard linear sector is beginning to
fail even in systems that are far beyond the traditional benchmark set.

This point is reinforced by the current state of practical rovibrational
workflows. Accurate quadratic force fields and semi-diagonal or partial cubic
force information are now becoming available also for medium-sized molecules,
and these data already support high-quality vibrational corrections to the
rotational constants. In that regime, a low-cost indicator of the limits of
the standard perturbative transport picture is particularly valuable. The
quartic $H_{22}$ diagnostic is exactly of this kind, because it is available
from the harmonic model alone, without requiring a full post-standard VPT4
treatment.

The same logic suggests a computationally adapted reading of the sextic
problem. The geometric sextic sector and the first sextic post-standard
diagnostic are already accessible from the harmonic model, while the explicit
cubic contribution may be further decomposed into a semi-diagonal part and a
genuinely three-index remainder. If the latter proves to be small in practice,
then a realistic hierarchy emerges for medium-sized systems: the standard
transport machinery, the quartic $H_{22}$ warning signal, and at least part of
the sextic structure may all be explored before a fully converged high-level
anharmonic treatment becomes affordable.

For water this hierarchy can already be tested explicitly. In the present
benchmark the genuinely three-index cubic remainder vanishes for all ten
standard sextic components, so that the full sextic constants are recovered
exactly from the geometry sector plus the semi-diagonal cubic contribution.
For the dominant $\Phi_{aaa}$ constant, for example, the geometry term is
about $7.72\times 10^5$ Hz and the semi-diagonal cubic contribution about
$9.08\times 10^4$ Hz, yielding the full value of about $8.62\times 10^5$ Hz
with no additional three-index correction. In this sense H$_2$O provides a
clean proof of principle for the computational hierarchy proposed here.

Formaldehyde provides a complementary and more realistic case. Here the
genuinely three-index cubic remainder is no longer zero, but it remains
subleading for the largest sextic constants. For the dominant
$\Phi_{aaa}$ contribution, for instance, the geometry term is about
$4.11\times 10^3$ Hz, the semi-diagonal cubic correction about
$-5.58\times 10^2$ Hz, and the genuine three-index remainder about
$1.18\times 10^2$ Hz, yielding a full value of about $3.67\times 10^3$ Hz.
The three-index part therefore amounts to only about $3\%$ of the total in
the leading channel, even though its relative effect becomes larger for
smaller components. Taken together, H$_2$O and H$_2$CO show that the
semi-diagonal cubic sector can already capture the dominant sextic structure,
while the genuinely three-index part acts as a systematic refinement rather
than as a qualitatively separate contribution.

The same trend persists for thioformaldehyde. For H$_2$CS the dominant
$\Phi_{aaa}$ contribution is about $5.05\times 10^3$ Hz at the geometric
level, with a semi-diagonal cubic correction of about $-3.76\times 10^1$ Hz
and a genuinely three-index remainder of only about $3.3$ Hz, corresponding
to less than $0.1\%$ of the total. Table~\ref{tab:sextic_cubic_hierarchy}
summarizes the dominant sextic component for H$_2$O, H$_2$CO, H$_2$CS, and
F$_2$CO, all reported after transport to a common $I^r$ representation.
F$_2$CO is especially informative because the genuinely three-index cubic
remainder is again very small in the leading channel, amounting to only about
$0.26\%$ of the total. Taken together, these examples suggest that the
genuinely three-index cubic sector is not universally negligible, but often
remains much smaller than the geometry term and the semi-diagonal cubic
contribution in the leading sextic channels. The calculations reported here
were all performed at the HPCS2 level; the purpose is therefore not a
benchmark-quality quantitative assessment of absolute sextic constants, but
the identification of robust trends in the hierarchy of geometric,
semi-diagonal cubic, and three-index cubic contributions.

\begin{table}[h]
\centering
\caption{Dominant sextic component and decomposition into geometry,
semi-diagonal cubic, and genuine three-index cubic contributions for four
illustrative molecules, all expressed in a common $I^r$ representation. All
values are in Hz. The calculations were performed at the HPCS2 level and are
intended to illustrate trends in the sextic hierarchy rather than
benchmark-quality absolute predictions.}

\begin{tabular}{lcccccc}
\hline
Molecule & Dominant component & Geometry & Cubic 2-index & Cubic 3-index & Total & 3-index / total \\
\hline
H$_2$O  & $\Phi_{aaa}$ & $7.72\times 10^5$ & $9.08\times 10^4$ & $0$ & $8.62\times 10^5$ & $0.0\%$ \\
H$_2$CO & $\Phi_{aaa}$ & $4.11\times 10^3$ & $-5.58\times 10^2$ & $1.18\times 10^2$ & $3.67\times 10^3$ & $3.2\%$ \\
H$_2$CS & $\Phi_{aaa}$ & $5.05\times 10^3$ & $-3.76\times 10^1$ & $3.32$ & $5.02\times 10^3$ & $0.07\%$ \\
F$_2$CO & $\Phi_{aaa}$ & $6.67\times 10^{-2}$ & $1.19\times 10^{-2}$ & $2.07\times 10^{-4}$ & $7.88\times 10^{-2}$ & $0.26\%$ \\
\hline
\end{tabular}

\label{tab:sextic_cubic_hierarchy}
\end{table}

This is not only a formal distinction. In finite-difference schemes based on
analytic gradients, the cost of the semi-diagonal cubic sector grows only
linearly with the number of vibrational modes, whereas the cost of the
genuinely three-index cubic sector grows quadratically. The separation of the
sextic cubic contribution into these two parts therefore mirrors a real
computational hierarchy, not just an algebraic rearrangement.\cite{cubic-vib_2024}

For this reason, the app and the corresponding Python implementation should be
seen as part of the scientific contribution rather than as auxiliary software.
They make the tensor transport and the associated diagnostics directly usable
on realistic spectroscopic datasets. In particular, the app implements the
standard quartic and sextic transport machinery discussed in this work and
also provides the first post-standard diagnostics, namely the quartic
$H_{22}$ term and its sextic linear analogue, from optional harmonic input
(`.fchk` or `xyz + Hessian`).

This practical scope is already illustrated in part by the HBN family
discussed above. Those systems are still small enough to admit careful PCS
treatment, but already large enough to make the low-cost hierarchy of the
present sextic formulation meaningful. The same workflow naturally extends to
larger flexible molecules, where repeated high-level refits quickly become
too expensive and a tensor-based prescreening of representations and
diagnostics becomes even more valuable.

\subsection{From elegant reformulation to spectroscopic tool}

Taken together, the examples above show why the present paper should not be
read as a purely formal note. DMSO shows that tensor transport is already
useful inside the standard sector because it separates true rovibrational
content from representation-induced numerical instability. H$_2$CO and
H$_2$CS show that the same logic extends cleanly to the standard sextic
sector. H$_2$O and H$_2$CO show that the first post-standard quartic
diagnostic, $H_{22}$, already behaves in the expected way on real molecules.

The resulting picture is therefore practical and hierarchical. Standard
quartics and standard sextics are transported linearly once lifted to their
tensor representatives. The first breakdown of that picture is measured by
$H_{22}$ on the quartic side and by the linear response in $\tau(H_{22})$ on
the sextic side. These quantities do not replace a full higher-order theory,
but they do provide exactly what is needed in day-to-day spectroscopic work:
representation transport, conditioning diagnostics, and early warning signs
that the standard perturbative picture is becoming insufficient.
